% ============== Section 7: AI-native physical layer (S43-S50) =============
\section{AI-Native Physical Layer}

% --- S43 -----------------------------------------------------------------
\begin{frame}{What ``AI-native'' actually means}
  \begin{itemize}\itemsep2pt
    \item 5G: AI bolted on (SON, scheduling). 6G ambition: AI as a
          \alert{design principle} inside the air interface
    \item \textbf{Model deficit}: no faithful math model (hardware
          nonlinearities, mmWave clutter) $\Rightarrow$ learn the model
    \item \textbf{Algorithm deficit}: model known, optimal algorithm
          intractable $\Rightarrow$ learn the algorithm
    \item Integration spectrum: replace a block $\to$ co-design blocks $\to$
          end-to-end learned PHY
  \end{itemize}
  \note{[S43, 1 min | clock 0:59:30, after Demo 4] The deficit taxonomy
  (Shlezinger/Eldar) is the most useful sorting hat: it predicts where
  learning helps and where classic DSP stays. Channel coding: no deficit,
  learning adds little. PA distortion compensation: model deficit, learning
  shines. Keep this frame for every AI-PHY paper you referee.}
\end{frame}

% --- S44 -----------------------------------------------------------------
\begin{frame}{Deep learning for channel estimation}
  \begin{itemize}\itemsep2pt
    \item Treat the pilot-grid estimate as a \alert{noisy image} $\Rightarrow$
          CNN denoising / super-resolution across time--frequency
    \item Learned estimators approach MMSE \emph{without} knowing channel
          statistics --- they absorb them from data
    \item Gains largest at low SNR and under model mismatch
    \item Costs: training data, generalization across cells/hardware,
          µs-latency inference budgets
  \end{itemize}
  \note{[S44, 1 min | clock 1:00:30] This is exactly Demo 5's setting ---
  hold the thought for ten minutes. The image analogy is not decoration: the
  TF channel grid has local correlation structure like natural images, which
  is why vision architectures transfer. The honest caveat list (data,
  generalization) returns on S49.}
\end{frame}

% --- S45 -----------------------------------------------------------------
\begin{frame}{Autoencoder end-to-end PHY}
  \begin{columns}[T]
    \column{0.55\textwidth}
    \begin{itemize}\itemsep2pt
      \item O'Shea \& Hoydis 2017: TX--channel--RX as one autoencoder; learn
            constellation and demapper \alert{jointly}
      \item Learned constellations rediscover --- and under impairments beat ---
            hand-designed ones
      \item Needs a differentiable channel $\Rightarrow$ surrogate models or RL
            for the real world
    \end{itemize}
    \column{0.43\textwidth}
    \centering
    \begin{tikzpicture}[node distance=2.5mm and 4mm]
      \node[blkb] (enc) {encoder\\NN};
      \node[blk, right=of enc, fill=cblue!15, draw=cblue] (ch) {channel};
      \node[blkb, right=of ch] (dec) {decoder\\NN};
      \draw[arr] (enc) -- (ch); \draw[arr] (ch) -- (dec);
      \node[lab, left=2mm of enc] (s) {$s$};
      \node[lab, right=2mm of dec] (sh) {$\hat s$};
      \draw[arr] (s) -- (enc); \draw[arr] (dec) -- (sh);
      \draw[arr, dashed, gray] (sh.south) to[bend left=25]
        node[lab, below]{end-to-end loss} (enc.south);
    \end{tikzpicture}
  \end{columns}
  \note{[S45, 1 min | clock 1:01:30] Historically the paper that launched the
  field: treat the whole link as one differentiable function and descend the
  gradient. Realism check for the audience: block-based systems will not
  vanish --- certification, interoperability and debugging all favor structure.
  Hybrids win in practice.}
\end{frame}

% --- S46 -----------------------------------------------------------------
\begin{frame}{Neural receivers}
  \begin{itemize}\itemsep2pt
    \item Replace equalization + demapping with one trained network over the
          received grid (channel estimate optional!)
    \item Demonstrated at 3GPP-compatible scale (multi-user MIMO neural
          receivers, Sionna-class toolchains)
    \item Robustness from training over channel \emph{distributions};
          complexity tamed by pruning/quantization
    \item Open operational question: retraining cadence vs channel drift
  \end{itemize}
  \note{[S46, 1 min | clock 1:02:30] The pragmatic middle ground that vendors
  actually prototype: keep the OFDM frame, swap the inner receiver. Notable
  detail: good neural receivers can skip explicit channel estimation ---
  pilots become just more input features. That should provoke the estimation
  theorists in the room --- in a good way.}
\end{frame}

% --- S47 -----------------------------------------------------------------
\begin{frame}{Model-based deep learning: deep unfolding}
  \begin{equation*}
    \text{ISTA: } \mathbf x^{(k+1)}=\eta_{\lambda}\!\left(\mathbf x^{(k)}
      +A^{H}(\mathbf y-A\mathbf x^{(k)})\right)
    \;\;\Rightarrow\;\;
    \text{LISTA: } \mathbf x^{(k+1)}=\eta_{\lambda_k}\!\left(W_1^{(k)}\mathbf y
      +W_2^{(k)}\mathbf x^{(k)}\right)
  \end{equation*}
  \begin{itemize}\itemsep2pt
    \item Unroll $K$ iterations of a classical algorithm into a $K$-layer
          network; \alert{learn the knobs} (steps, thresholds, matrices)
    \item Order-of-magnitude fewer iterations at equal accuracy; retains
          interpretability and per-layer guarantees flavor
    \item PHY wins: unfolded detection, beamforming, sparse channel estimation
  \end{itemize}
  \note{[S47, 1.5 min | clock 1:03:30] The best of both worlds slide. Sell it
  as ``your favorite iterative algorithm, after weight training'': same
  skeleton, data-tuned parameters, so priors and structure survive. For the
  faculty: unfolding is a paper machine --- take any classic iteration in your
  drawer and unfold it. This equation is the cheat-sheet's centerfold.}
\end{frame}

% --- S48 -----------------------------------------------------------------
\begin{frame}{Semantic communications}
  \begin{itemize}\itemsep2pt
    \item Shift the target from \alert{bit fidelity} to \alert{task fidelity}:
          transmit what the receiver needs to act, not every bit
    \item Neural joint source--channel coding: graceful degradation, no cliff
          effect --- compelling image/speech demos
    \item Open theory: a working semantic information measure; where the
          separation theorem stands; standardizable interfaces
    \item Fair positioning: promising for constrained links (XR, satellite
          IoT); not a 6G day-one feature
  \end{itemize}
  \note{[S48, 1 min | clock 1:05] Flex slide --- can merge into S49 if behind
  schedule. One vivid example: transmitting a face for a video call needs the
  expression, not the pixels. The cliff-effect contrast: digital codecs die
  abruptly at threshold SNR; learned JSCC degrades like analog TV --- soft.}
\end{frame}

% --- S49 -----------------------------------------------------------------
\begin{frame}{The hard problems: data, generalization, trust}
  \begin{itemize}\itemsep2pt
    \item \textbf{Data}: channel measurements are vendor-siloed; sim-to-real
          transfer is the quiet blocker
    \item \textbf{Generalization}: across cells, hardware, mobility;
          distribution-shift monitoring in a live network
    \item \textbf{Budget}: µs-scale inference, joules per inference, PHY
          real-time constraints
    \item \textbf{Trust}: testability, certification, explainability ---
          telecom-grade requirements
    \item Community opening: standardized channel datasets \& benchmarks would
          be a genuine service to the field
  \end{itemize}
  \note{[S49, 1 min | clock 1:06] The referee-hat slide. Encourage the room:
  a well-curated public dataset + benchmark for one PHY task can out-impact
  another architecture paper. If S48 was skipped, borrow its minute here for
  discussion of sim-to-real.}
\end{frame}

% --- S50 -----------------------------------------------------------------
\begin{frame}{AI/ML in 3GPP}
  \begin{itemize}\itemsep1pt
    \item Rel-18 study $\to$ Rel-19 normative first steps; three use cases:
          \alert{CSI feedback compression}, \alert{beam prediction},
          \alert{positioning}
    \item Two-sided models (UE encoder / network decoder): the
          interoperability puzzle of the decade
    \item 6G study (Rel-20): AI-native candidates; the real spec work is
          \emph{life-cycle management} --- monitoring, fallback, retraining
  \end{itemize}
  \demoslot{5 (2.5 min)}{\textbf{LS vs MMSE vs a small neural network} for OFDM
    channel estimation --- NMSE curves, pre-trained net, runs in seconds.
    \hfill\src{matlab/demo5\_dl\_channel\_estimation}}
  \note{[S50 + DEMO 5, 1+2.5 min | clock 1:07] Standards reality in three
  bullets, then the last MATLAB switch. Demo 5 narration: LS is the
  statistics-free baseline; MMSE knows the true covariance (unfair advantage);
  the small NN learns from data and lands within a fraction of a dB of MMSE.
  That is the whole AI-PHY promise in one plot. TIMING NOTE: this demo is the
  first to sacrifice if late --- appendix A5 shows the same curves. Done by
  ~1:10.}
\end{frame}
